\section{Conclusiones}

Las conclusiones se presentan en el mismo orden que los seis objetivos
específicos definidos en el Capítulo 1 y se sustentan en las evidencias
de análisis, diseño, implementación y pruebas del sistema.

\begin{enumerate}
    \item El primer objetivo específico se cumplió mediante la
    especificación de los requisitos funcionales y no funcionales,
    la identificación de actores y la delimitación de las operaciones
    prioritarias. La evidencia se encuentra en el análisis de requisitos
    y en el diagrama general de casos de uso, que relacionan la gestión
    de pacientes, las solicitudes radiográficas y el acceso temporal
    utilizado por CERPAX.

    \item El segundo objetivo específico se cumplió para el componente
    web mediante el diseño de la arquitectura, el modelo de datos,
    las interfaces y el contrato de la API. Los diagramas de arquitectura
    y datos, el diccionario de la entidad principal y las capturas de
    las pantallas respaldan este resultado. El cumplimiento del objetivo
    conjunto web y móvil es parcial, porque las evidencias revisadas
    no incluyen la verificación del componente móvil.

    \item El tercer objetivo específico se cumplió para las
    funcionalidades web priorizadas mediante la implementación de
    operaciones de registro, consulta, actualización y eliminación de
    pacientes, asignación a odontólogos, creación de solicitudes
    radiográficas y generación de enlaces temporales y códigos QR.
    Las pruebas locales también verificaron la carga multipart,
    la asociación del archivo con la solicitud y su consulta autorizada.
    El código fuente y los casos de prueba documentados respaldan
    estas funcionalidades; la prueba local de archivos no acredita
    la entrega directa a Supabase Storage en producción.

    \item El cuarto objetivo específico se cumplió en los controles
    web evaluados mediante autenticación con contraseñas protegidas
    con bcrypt, sesiones JWT, autorización en el servidor y accesos
    temporales vinculados a solicitudes. Las pruebas comprobaron
    el rechazo de credenciales incorrectas, sesiones inválidas,
    consultas administrativas realizadas por odontólogos y accesos
    a pacientes no asignados. También confirmaron que el paciente
    no puede iniciar sesión y que los enlaces revocados o vencidos
    son rechazados.

    \item El quinto objetivo específico se cumplió para los requisitos
    evaluados mediante pruebas reproducibles del backend y un recorrido
    de navegador. La ejecución final del backend registró 31 pruebas
    aprobadas, ninguna fallida y ninguna omitida. Durante la verificación
    se detectó y corrigió la carga anticipada de la pantalla de
    radiografías, y el recorrido posterior fue aprobado. Asimismo,
    las diez consultas de pacientes medidas en Vercel respondieron
    entre 208 y 349 milisegundos, cumpliendo RNF-01 en esa muestra.
    Los reportes y las evidencias de la Sección 2.8 respaldan estos
    resultados, cuyo alcance corresponde a las condiciones descritas.

    \item El sexto objetivo específico se cumplió parcialmente:
    la plataforma web dispone de una dirección pública en Vercel,
    un repositorio con control de versiones y documentación de
    configuración y despliegue. Las comprobaciones documentadas
    confirmaron la disponibilidad de la API, el inicio de sesión
    y la consulta de solicitudes radiográficas. Sin embargo,
    la revisión local que incorpora las últimas correcciones
    requiere publicación y verificación, y la entrega completa
    a Supabase Storage desde el entorno de producción no fue
    validada durante esta revisión.
\end{enumerate}
